% ============================================================================
% Fichier  : partie6/P06_C20_mesurer_recherche.tex
% Titre    : Mesurer la Supervision et la Réponse
% Auteur   : MINKA MI NGUIDJOÏ Thierry Emmanuel
% ============================================================================
% Sources : ISO/IEC 27004 ; cadre OUI/NON/NEK de l'auditologie (Minka). MBAM.
% ============================================================================

\chapter{Mesurer la Supervision et la Réponse}
\label{chap:mesurer}
\chaptermark{Mesurer et conclure}

\epigraph{%
    « Mesurer la sécurité, c'est aussi savoir dire ce que l'on ignore. »%
}{Principe du cours}

\niveauChapitre{Avancé}{L'ensemble du cours~; notions d'auditologie.}

\begin{boxobjectifs}
\begin{itemize}
    \item Choisir des indicateurs de supervision et de réponse honnêtes.
    \item Appliquer le verdict à trois valeurs OUI, NON, NEK.
    \item Éviter les indicateurs trompeurs.
    \item Inscrire la supervision dans un cycle d'amélioration.
    \item Ouvrir sur la recherche en détection et réponse.
\end{itemize}
\end{boxobjectifs}

\bigskip

% ============================================================================
\section*{Question directrice : LiZbiZ est-elle « prête » pour le prochain incident~?}
% ============================================================================

Six mois après, la question revient~: « Sommes-nous prêts~? » La réponse
honnête n'est ni oui ni non en bloc, mais un ensemble de verdicts précis, et,
pour ce qui n'a pu être éprouvé, un franc NEK. Ce chapitre clôt le cours sur
la mesure.

\medskip
\textbf{Plan du chapitre.}
\begin{enumerate}
    \item Des indicateurs honnêtes (section~\ref{sec:c20-indicateurs}).
    \item Le verdict à trois valeurs (section~\ref{sec:c20-verdict}).
    \item Les indicateurs trompeurs (section~\ref{sec:c20-trompeurs}).
    \item Le cycle et la recherche (section~\ref{sec:c20-cycle}).
\end{enumerate}

\separateur

% ============================================================================
\section{Des Indicateurs Honnêtes}
\label{sec:c20-indicateurs}
% ============================================================================

\begin{tableaucours}{L{4.8cm}L{7.4cm}}{%
    \tableauHeader{Indicateur} &
    \tableauHeader{Ce qu'il dit}
}
Délai de détection, de réaction & Rapidité à voir et à contenir \\
Couverture de journalisation & Part du parc réellement surveillé \\
Couverture de détection (ATT\&CK) & Part des techniques que l'on verrait \\
Taux de réussite des exercices & Préparation de la cellule \\
Taux de correction des actions RETEX & Capacité à apprendre \\
\end{tableaucours}
\vspace{-0.8em}
\begin{center}{\small\itshape Tableau~20.1. Indicateurs de supervision et de réponse~\cite{iso27004}}\end{center}

% ============================================================================
\section{Le Verdict à Trois Valeurs}
\label{sec:c20-verdict}
% ============================================================================

\begin{boxdef}
\termecle{Verdict à trois valeurs}~: à chaque question de préparation, la
réponse est \sigle{OUI} (démontré, avec preuve), \sigle{NON} (réfuté, avec
preuve) ou \sigle{NEK} (indéterminé~: les preuves manquent pour trancher). Le
NEK n'est pas un aveu de faiblesse~: c'est une exigence d'honnêteté qui
désigne où porter l'effort.
\end{boxdef}

\begin{tableaucours}{L{2.0cm}L{4.8cm}L{5.4cm}}{%
    \tableauHeader{Verdict} &
    \tableauHeader{Question} &
    \tableauHeader{Exemple (LiZbiZ)}
}
OUI & Détecte-t-on un compte de service anormal~? & Règle écrite et éprouvée \\
NON & Détecte-t-on l'exfiltration depuis tout serveur~? & FS-01 journalisé, pas les autres \\
NEK & Résisterait-on à un attaquant plus patient~? & Jamais éprouvé par un exercice \\
\end{tableaucours}
\vspace{-0.8em}
\begin{center}{\small\itshape Tableau~20.2. Le verdict OUI / NON / NEK appliqué~\cite{minka2026auditologie}}\end{center}

\begin{boxalert}
\textbf{Violation typique.} Conclure OUI (« nous sommes prêts ») faute d'avoir
éprouvé.

\smallskip
\textbf{Règle~:} ce qui n'a pas été testé par un exercice reste NEK, pas OUI.
L'absence d'incident n'est pas une preuve de préparation.
\end{boxalert}

% ============================================================================
\section{Les Indicateurs Trompeurs}
\label{sec:c20-trompeurs}
% ============================================================================

\begin{tableaucours}{L{4.4cm}L{7.8cm}}{%
    \tableauHeader{Indicateur trompeur} &
    \tableauHeader{Pourquoi il trompe}
}
Nombre d'alertes & Dépend des règles, pas du risque \\
« Aucun incident ce mois-ci » & Peut signifier « rien détecté » \\
Nombre d'outils déployés & L'outil n'est pas la détection \\
Temps passé au SOC & Mesure l'effort, pas le résultat \\
\end{tableaucours}
\vspace{-0.8em}
\begin{center}{\small\itshape Tableau~20.3. Indicateurs à manier avec prudence}\end{center}

% ============================================================================
\section{Le Cycle et la Recherche}
\label{sec:c20-cycle}
% ============================================================================

La supervision s'inscrit dans un cycle~: détecter, répondre, apprendre,
améliorer la détection, éprouver par des exercices, repartir des NEK. La
recherche prolonge ce cycle~: mesure de la couverture de détection,
automatisation de la réponse, renseignement adapté au contexte
d'Afrique centrale, et application du cadre de l'auditologie à la réduction
de l'incertitude en supervision.

\begin{boxscenario}{Opération MBAM. L'exercice qui lève des NEK}
\textbf{Constatation.} Pour éprouver sa préparation, LiZbiZ participe à un
exercice d'équipe (type équipe violette) avec NYONG SÉCURITÉ~: 30~techniques
jouées, 18~détectées. Les 12~non détectées étaient des NEK~: on ne savait pas
si on les verrait. Désormais, on sait que non, et on écrit les règles.

\medskip
\textbf{Analyse} (confiance élevée). L'exercice transforme des NEK en OUI ou
en NON étayés. C'est le seul moyen honnête de répondre « sommes-nous
prêts~? ».
\end{boxscenario}

\begin{boxmetier}{RSSI, analyste SOC, analyste forensique}
\textbf{Dire le vrai, NEK compris.} Le professionnel ne cède pas au OUI
rassurant~: il distingue ce qu'il a démontré, réfuté, et ce qu'il n'a pu
trancher. C'est à ce triple langage qu'on reconnaît un bilan honnête.
\end{boxmetier}

\separateur

\begin{boxresume}
\textbf{Ce qu'il faut retenir de ce chapitre~:}
\begin{itemize}
    \item \textbf{Des indicateurs} de délai, couverture, exercices,
          correction.
    \item \textbf{OUI, NON, NEK}~: le verdict honnête.
    \item \textbf{L'absence d'incident} n'est pas une preuve de préparation.
    \item \textbf{Les indicateurs d'activité} trompent.
    \item \textbf{Les exercices} transforment les NEK en verdicts étayés.
\end{itemize}
\end{boxresume}

\begin{boxinfo}
\textbf{Points de contrôle de l'élève.} Je sais~: choisir des indicateurs
honnêtes~; appliquer OUI/NON/NEK~; reconnaître un indicateur trompeur~;
décrire le cycle~; situer les pistes de recherche.
\end{boxinfo}

\bigskip

\begin{boxexo}[title={\ding{45}\quad Exercice~20.1. Indicateurs \niveaubase}]
Citez trois indicateurs honnêtes de supervision.
\end{boxexo}

\begin{boxexo}[title={\ding{45}\quad Exercice~20.2. Trois valeurs \niveaubase}]
Donnez le sens de OUI, NON et NEK.
\end{boxexo}

\begin{boxexo}[title={\ding{45}\quad Exercice~20.3. Trompeur \niveaubase}]
Pourquoi « aucun incident ce mois-ci » peut-il tromper~?
\end{boxexo}

\begin{boxexo}[title={\ding{45}\quad Exercice~20.4. Classer \niveaumoyen}]
Attribuez OUI, NON ou NEK~: détection d'un compte de service anormal~;
résistance à un attaquant patient~; détection de l'exfiltration partout.
\end{boxexo}

\begin{boxexo}[title={\ding{45}\quad Exercice~20.5. Pas de test \niveaumoyen}]
Pourquoi ce qui n'a pas été éprouvé reste-t-il NEK~?
\end{boxexo}

\begin{boxexo}[title={\ding{45}\quad Exercice~20.6. MBAM \niveaumoyen}]
Comment l'exercice MBAM transforme-t-il des NEK en verdicts étayés~?
\end{boxexo}

\begin{boxexo}[title={\ding{45}\quad Exercice~20.7. Note au DG \niveaumoyen}]
Répondez à « sommes-nous prêts~? » en OUI / NON / NEK pour LiZbiZ.
\end{boxexo}

\begin{boxexo}[title={\ding{45}\quad Exercice~20.8. Outils \niveaumoyen}]
Pourquoi le nombre d'outils déployés n'est-il pas un bon indicateur~?
\end{boxexo}

\begin{boxexo}[title={\ding{45}\quad Exercice~20.9. Tableau de bord \niveauexpert}]
Construisez le tableau de bord annuel de supervision et de réponse de LiZbiZ,
NEK compris.
\end{boxexo}

\begin{boxexo}[title={\ding{45}\quad Exercice~20.10. Question de recherche \niveauexpert}]
Formulez une question de recherche sur la mesure de la couverture de
détection en contexte d'Afrique centrale.
\end{boxexo}

\bigskip

\begin{boxinfo}
\textbf{Pour aller plus loin.}
\begin{itemize}
    \item ISO/IEC 27004~\cite{iso27004}~; cadre de
          l'auditologie~\cite{minka2026auditologie}.
\end{itemize}
\end{boxinfo}

\begin{boxinfo}
\textbf{Fin du cours.} De l'écran de rançon au verdict mesuré, SIS589 aura
montré qu'une attaque se voit venir, s'investigue proprement et se maîtrise~:
superviser, investiguer, répondre, apprendre, mesurer, sans jamais confondre
ce que l'on constate et ce que l'on en déduit.
\end{boxinfo}
